\subsection{Midpoint of a Segment} 
The \emph{midpoint} $M$ of a segment $AB$ is the point on $\overline{AB}$ such that $AM=MB$. This is easy to find if $\overline{AB}$ is \makeblank[1.5in] or \makeblank[1.5in].
\begin{example}
Let $A=(-3,2)$, $B=(-3,0)$, $C=(-1,1)$, and $D=(3,1)$. Plot these points, and find the midpoints  of $\overline{AB}$ and $\overline{CD}$.
\begin{lpic}{}{4}
\foreach \y/\lbl in {-1/AB,-3/CD} \regtextleft{1}{\y}{$\overline{\lbl}$ has midpoint};
\end{lpic}
\end{example}
In general:
\begin{itemize}
\item If $A=(x,y_A)$ and $B=(x,y_B)$ (same \makeblanksmall coordinate), the midpoint of $\overline{AB}$ is $M=\left(\makeblanksmall,\makeblank[1in]\right)$
\item If $A=(x_A,y)$ and $B=(x_B,y)$ (same \makeblanksmall coordinate), the midpoint of $\overline{AB}$ is $M=\left(\makeblank[1in],\makeblanksmall\right)$
\end{itemize}
For other segments, we use the properties of congruent triangles.
\begin{example}
Let $A=(-2,-3)$ and $B=(4,1)$. Plot these points, and find the midpoint of $\overline{AB}$.
\begin{lpic}{}{6}
\end{lpic}
\end{example}
\begin{displaybox}[Midpoint Formula]
If $A=(x_A,y_A)$ and $B=(x_B,y_B)$, then the midpoint of $\overline{AB}$ is
\[
M=\left(\frac{x_A+x_B}{2},\frac{y_A+y_B}{2}\right)
\]
\end{displaybox}